% Generated draft: requires mathematical and source review.
\documentclass[11pt]{article}
\usepackage[margin=25mm]{geometry}
\usepackage{fontspec}
\setmainfont{Latin Modern Roman}
\usepackage{amsmath,amssymb,mathtools,unicode-math}
\setmathfont{Latin Modern Math}
\usepackage{enumitem,needspace,xurl,hyperref}
\setlength{\parindent}{0pt}
\setlength{\parskip}{5pt}
\setlength{\emergencystretch}{3em}
\newcommand{\N}{\mathbb{N}}
\newcommand{\Z}{\mathbb{Z}}
\newcommand{\Q}{\mathbb{Q}}
\newcommand{\R}{\mathbb{R}}
\newcommand{\C}{\mathbb{C}}
\newcommand{\F}{\mathbb{F}}
\newcommand{\A}{\mathbb{A}}
\newcommand{\PP}{\mathbb{P}}
\newcommand{\E}{\mathbb{E}}
\newcommand{\eps}{\varepsilon}
\newcommand{\dd}{\,\mathrm{d}}
\begin{document}
\section*{UnsolvedMath ID: 30005283}
\textbf{Source problem code:} catalogue entry
\subsection{Definitions and mathematical statement}
Let \(G\) be a definable group of finite Morley rank acting definably on a definable set \(X\) in a common ambient structure. Write \(\operatorname{RM}(Y)\) for the Morley rank of a definable set \(Y\), with \(\operatorname{RM}(\varnothing)=-\infty\). The supplied finite-Morley-rank setting implicitly requires \(r:=\operatorname{RM}(X)\) to be finite. For a positive integer \(k\), the diagonal action on \(X^k\) is
\[
 g\cdot(x_1,\ldots,x_k)=(g\cdot x_1,\ldots,g\cdot x_k).
\]
The action is \emph{generically \(k\)-transitive} if there is an orbit \(O\subseteq X^k\) such that
\[
 \operatorname{RM}(X^k\setminus O)<\operatorname{RM}(X^k).
\]
This is the definition in the supplied excerpt; it uses a complement of strictly smaller rank, rather than merely an orbit of full rank.

The supplied record attributes to Borovik and Cherlin the conjectural bound and rigidity assertion
\[
 k\leq r+2,
\]
with equality only if, for some integer \(n\) and a suitable field \(K\), the action is definably isomorphic to the natural action
\[
 \operatorname{PGL}_n(K)\curvearrowright\mathbb P^{n-1}(K).
\]
Here \(\operatorname{PGL}_n(K)=\operatorname{GL}_n(K)/\{aI_n:a\in K^\times\}\), and \(\mathbb P^{n-1}(K)\) is the set of one-dimensional \(K\)-linear subspaces of \(K^n\). A definable isomorphism of actions means a definable group isomorphism \(\phi:G\to\operatorname{PGL}_n(K)\) and a definable bijection \(f:X\to\mathbb P^{n-1}(K)\) satisfying
\[
 f(g\cdot x)=\phi(g)\cdot f(x)
 \qquad(g\in G,\ x\in X).
\]

\emph{Unresolved qualifications.} The supplied formulation does not specify positive rank or infinitude of \(X\), faithfulness of the action, connectedness or transitivity assumptions, the field \(K\), its ambient definable structure, or the range of \(n\). Rank-zero degeneracies obstruct the unrestricted bound, and a definable kernel acting trivially obstructs rigidity for the whole group as stated. It is therefore necessary to determine whether faithfulness is intended or whether rigidity concerns the induced faithful group \(G/\ker(G\curvearrowright X)\), where
\[
 \ker(G\curvearrowright X)=\{g\in G:(\forall x\in X)\ g\cdot x=x\}.
\]
In the customary projective interpretation with an algebraically closed field \(K\) of Morley rank one, one would have \(r=n-1\) and hence equality would mean \(k=n+1\); those field hypotheses are not supplied explicitly. The displayed assertion records the source's intended target, without selecting missing hypotheses or claiming a complete corrected conjecture.
\subsection{Short English statement}
The supplied Borovik--Cherlin conjecture asks whether generic transitivity is bounded by the Morley rank of the acted-on set plus two, with equality forcing the standard projective linear action. The source leaves essential hypotheses unresolved, especially nondegeneracy, faithfulness, and the field in the equality case. This is a statement-only draft requiring mathematical review.
\subsection{Sources}
\begin{itemize}
\item Local UnsolvedMath registry: \texttt{data/unsolvedmath/problems.json}, permanent ID \(30005283\), title ``Borovik--Cherlin Generic-Transitivity Conjecture,'' exact source problem code \texttt{catalogue entry}. The supplied record attributes the conjecture to Borovik and Cherlin; it gives no individual proposer.
\item Supplied report citation: \emph{Mini-Workshop: Topological and Differential Expansions of o-minimal Structures} (2023), Oberwolfach Reports, pp.~3005--3050, DOI \texttt{10.4171/owr/2022/52}. The excerpt carries page number 3029. \url{https://doi.org/10.4171/owr/2022/52}. The registry separately gives proposed year 2022.
\item Supplied literature reference: Borovik--Nesin--Wagner, \emph{Bounding generic sharp transitivity}. \url{https://arxiv.org/abs/2407.09636}.
\item Supplied literature reference: Borovik--Burdges, \emph{J. London Math. Soc.} \url{https://doi.org/10.1112/jlms/jdm088}. No article title is supplied.
\item Supplied literature reference: EMS Press, OWR 2022/52. \url{https://doi.org/10.4171/owr/2022/52}.
\item Literature assessment dated exactly \texttt{2026-08-22}: the supplied curation labels the problem ``partially solved'' and reports important special cases, including algebraic-group settings, with continuing partial results for the general finite-Morley-rank statement. This is supplied dated curation, not independently verified literature assessment. The registry also labels the statement \texttt{corrected\_verified} and reports checking against the official report; those are supplied provenance claims, not verification performed for this draft.
\end{itemize}
\end{document}
